\chapter{Collaudo e Testing}

La fase di verifica e validazione del sistema \textbf{EscapeManager} è stata condotta seguendo un approccio multilivello, utilizzando il framework \textbf{JUnit 5}. L'obiettivo del collaudo non è stato solo quello di garantire l'assenza di difetti evidenti (bug), ma soprattutto di dimostrare la robustezza dell'architettura e la perfetta aderenza ai requisiti funzionali delineati nel Capitolo 2.

\section{Strategia di Verifica}
La strategia di test è organizzata su tre livelli:
\begin{itemize}
    \item \textbf{Unit Test (White-Box):} verificano le logiche pure del Domain Model e dei Design Pattern.
    \item \textbf{Integration Test (DAO/DB):} verificano il mapping ORM e la correttezza delle query principali.
    \item \textbf{Functional Test (Black-Box):} esercitano i Controller come surrogati della UI, in linea con i template UC.
\end{itemize}

\section{Setup e Dati di Test}
Per garantire test ripetibili, prima di ogni test di integrazione o funzionale il database viene riportato ad uno stato noto tramite \texttt{DatabaseHelper.resetTestDatabase()}, che esegue TRUNCATE e reinserisce dati minimi consistenti (una stanza e due clienti).

\section{Testing Strutturale (White-Box)}

I test strutturali (o Unit Test) si sono concentrati sulle classi del \texttt{Domain Model}, con l'obiettivo di massimizzare la copertura (Code Coverage) dei rami condizionali e verificare la correttezza algoritmica, in totale isolamento dalle dipendenze esterne.

Particolare attenzione è stata dedicata al collaudo dei Design Pattern. Il Listing \ref{lst:test_strategy} mostra un test parametrizzato per la verifica delle politiche di prezzo dinamiche (\textit{Strategy Pattern}).

\begin{lstlisting}[caption={Unit Test per la logica di calcolo del prezzo (Strategy Pattern)}, label={lst:test_strategy}]
class PricingStrategyTest {

    @Test
    void testTariffaBaseCalcoloCorretto() {
        // Setup
        Stanza stanza = new Stanza("R01", "FI01", "Horror Asylum", 6, 20.0);
        stanza.setPricingStrategy(new TariffaBase());
        
        // Execution: 4 giocatori a 20 euro base = 80 euro
        double prezzoCalcolato = stanza.calcolaPrezzo(4);
        
        // Assertion
        assertEquals(80.0, prezzoCalcolato);
    }
}
\end{lstlisting}

Oltre al calcolo del prezzo, i test unitari verificano anche le transizioni di stato della \texttt{Stanza} (State Pattern) e la corretta gestione delle eccezioni in caso di transizioni non valide.

\section{Testing di Integrazione (DAO e Database)}

Per verificare il corretto interfacciamento tra l'applicazione Java e il database PostgreSQL (livello ORM), sono stati implementati test di integrazione sulle classi DAO.
Per garantire la \textbf{ripetibilità} e l'isolamento, il costrutto \texttt{@BeforeEach} invoca \texttt{DatabaseHelper.resetTestDatabase()}, che esegue TRUNCATE e reinserisce un dataset minimo prima di ogni metodo di test.

\begin{lstlisting}[caption={Integration Test per StanzaDAOPostgres con reset dello stato}, label={lst:test_dao}]
class StanzaDAOPostgresTest {
    private StanzaDAO stanzaDao;

    @BeforeEach
    void setUp() {
        DatabaseHelper.resetTestDatabase();
        stanzaDao = DAOFactory.getDAOFactory().getStanzaDAO();
    }

    @Test
    void testSaveAndFindById() {
        Stanza stanza = stanzaDao.findById("R01");
        assertNotNull(stanza);
        assertEquals("Horror Asylum", stanza.getTema());
    }
}
\end{lstlisting}

\section{Testing Funzionale (Black-Box e Tracciabilità)}

Il livello più alto di collaudo ha riguardato i \texttt{Controllers}. In questa fase (Black-Box Testing) la logica interna viene ignorata e il sistema viene testato esclusivamente in base agli input e agli output attesi.

Per garantire la massima \textbf{tracciabilità}, ogni test funzionale è stato mappato direttamente su uno specifico \textit{Use Case Template} (Sezione 2.2). Sono stati collaudati sia i \textit{Flussi Base} che i \textit{Flussi Alternativi} (le eccezioni).

Il Listing \ref{lst:test_funzionale} dimostra come il sistema gestisce fluidamente la \textit{Race Condition} descritta nell'UC1 (Alternativa 3a), negando la prenotazione se lo slot risulta già occupato a database.

\begin{lstlisting}[caption={Functional Test tracciato sull'UC1 (Flusso Alternativo 3a)}, label={lst:test_funzionale}]
class PrenotazioneControllerTest {
    private PrenotazioneController controller;
    private PrenotazioneDAO prenotazioneDAO;

    @BeforeEach
    void setUp() {
        DatabaseHelper.resetTestDatabase();
        DAOFactory factory = DAOFactory.getDAOFactory();
        prenotazioneDAO = factory.getPrenotazioneDAO();
        StanzaDAO stanzaDAO = factory.getStanzaDAO();
        controller = new PrenotazioneController(prenotazioneDAO, stanzaDAO);
    }

    @Test
    void testPrenotazioneNegataPerSlotOccupato_UC1_Alt3a() {
        // Setup: Il Cliente 2 tenta di prenotare uno slot appena occupato dal Cliente 1
        String stanzaId = "R01";
        LocalDateTime slotRichiesto = LocalDateTime.of(2026, 10, 31, 21, 0);
        
        // Simulo la race condition pre-occupando lo slot via DAO
        Prenotazione esistente = new Prenotazione(
            UUID.randomUUID().toString(), "CLIENTE_01", stanzaId,
            slotRichiesto, 3, 60.0, "CONFERMATA");
        prenotazioneDAO.save(esistente);

        // Execution & Assertion
        SlotNonDisponibileException eccezione = assertThrows(
            SlotNonDisponibileException.class, 
            () -> controller.effettuaPrenotazione("CLIENTE_02", stanzaId, slotRichiesto, 4)
        );
        
        assertEquals("Lo slot orario selezionato non e' piu' disponibile.", eccezione.getMessage());
    }
}
\end{lstlisting}

\subsection{Catalogo dei Test}
La Tabella \ref{tab:test_catalog} riassume l'intera suite di test, indicando tipo e identificativo.

\begin{table}[H]
    \centering
    \small
    \renewcommand{\arraystretch}{1.2}
    \begin{tabularx}{\textwidth}{|>{\bfseries}l|>{\bfseries}l|X|}
        \hline
        ID & Tipo & Test \\ \hline
        U1 & Unit & \texttt{PricingStrategyTest.testTariffaBaseCalcoloCorretto} \\ \hline
        U2 & Unit & \texttt{StanzaTest.testTransizioneCorretta} \\ \hline
        U3 & Unit & \texttt{StanzaTest.testEccezioneStatoNonValido} \\ \hline
        I1 & Integration & \texttt{StanzaDAOPostgresTest.testSaveAndFindById} \\ \hline
        F1 & Functional & \texttt{PrenotazioneControllerTest.testPrenotazioneNegataPerSlotOccupato\_UC1\_Alt3a} \\ \hline
        F2 & Functional & \texttt{ListaAttesaControllerTest.testIscrizioneListaAttesa\_OK} \\ \hline
        F3 & Functional & \texttt{ListaAttesaControllerTest.testIscrizioneListaAttesa\_Duplicato} \\ \hline
        F4 & Functional & \texttt{GestioneStatoControllerTest.testEseguiTransizioneAvvia} \\ \hline
        F5 & Functional & \texttt{GestioneStatoControllerTest.testTransizioneNonValida\_TerminaSuDisponibile\_UC3\_Alt4a} \\ \hline
        F6 & Functional & \texttt{TariffeControllerTest.testCambioTariffaWeekend\_UC4} \\ \hline
    \end{tabularx}
    \caption{Catalogo dei test (unit, integrazione, funzionali)}
    \label{tab:test_catalog}
\end{table}

\subsection{Tracciabilità Use Case \textrightarrow{} Test}
I test funzionali F1--F6 sono tracciati direttamente sui template del Capitolo 2. La Tabella \ref{tab:traceability} sintetizza la copertura dei flussi base e alternativi.

\begin{table}[H]
    \centering
    \small
    \renewcommand{\arraystretch}{1.2}
    \begin{tabularx}{\textwidth}{|>{\bfseries}l|X|}
        \hline
        Use Case / Flusso & Test ID \\ \hline
        UC1 - Flusso alternativo 3a (slot occupato) & F1 \\ \hline
        UC2 - Flusso base (iscrizione) & F2 \\ \hline
        UC2 - Alternativa 5a (duplicato) & F3 \\ \hline
        UC3 - Flusso base (avvio sessione) & F4 \\ \hline
        UC3 - Alternativa 4a (transizione non valida) & F5 \\ \hline
        UC4 - Flusso base (cambio tariffa) & F6 \\ \hline
    \end{tabularx}
    \caption{Matrice di tracciabilità Use Case \textrightarrow{} Test}
    \label{tab:traceability}
\end{table}

\section{Risultati}
La suite di test è composta da \textbf{7 classi} e \textbf{10 test cases}, distribuiti tra unit, integrazione e test funzionali. Non è stata eseguita una misura formale di coverage, ma i test coprono i principali rami decisionali dei pattern e i flussi alternativi più critici nei casi d'uso.

\section{Conclusioni}

L'approccio metodologico adottato, guidato dai principi dell'Object-Oriented Design e supportato dall'uso estensivo dei Design Pattern, ha permesso di sviluppare \textbf{EscapeManager} come un sistema robusto, altamente coeso e disaccoppiato. 
La suite di collaudo ha confermato la solidità della logica di business e la corretta integrazione con il layer relazionale, soddisfacendo integralmente i requisiti posti in fase di analisi e fornendo un'architettura facilmente estensibile per sviluppi futuri (es. integrazione di sensori fisici nelle stanze tramite \textit{Adapter Pattern} o implementazione di un'interfaccia grafica avanzata al posto dell'attuale CLI).


